% ==============================================================================================
\section{Protocols and message formats}
\label{sec:proto}

This section describes how the programs talk: the paths, the messages, the timers, the retries and how duplicates are handled. The full tables are in \texttt{docs/protocol.md} and \texttt{docs/shared-log.md} \cite{repo}.

\subsection{The life of one report}

The easiest way to see the protocol is to follow one real report. Figure~\ref{fig:seq-report} draws one report from a run at 10,000 reports a second, taken line by line from the event logs of all four PCs it touched (\texttt{analysis/results/request-trace-example.txt}). The times count from the client's send. They come from four different clocks, which were within a few milliseconds of each other, so they are close but not exact.

\begin{figure}[!htb]
\centering
\begin{tikzpicture}[x=1cm,y=1cm]
  \seqactor[clientorange!12]{0}{Client 3\\(MeshClient)}{-10.6}
  \seqactor{4.4}{Node 1}{-10.6}
  \seqactor{8.8}{Node 2}{-10.6}
  \seqactor{13.2}{Node 3}{-10.6}
  % the group send
  \seqmsg{0}{13.2}{-1.0}{}
  \seqmsg{0}{8.8}{-1.25}{}
  \seqmsg{0}{4.4}{-1.5}{}
  \node[mlabel] at (6.6,-0.95) {0 ms: the same report to all 3 nodes, UDP};
  % bids
  \seqmsg{13.2}{4.4}{-2.2}{11.8 ms: node 3 bids 49}
  \seqmsg{8.8}{13.2}{-2.75}{12.7 ms: node 2 bids 135}
  \seqmsg{4.4}{13.2}{-3.3}{26.7 ms: node 1 bids 123}
  \seqnote{2.2}{-2.75}{3.4}{each bid goes to both peers;\\ one arrow is drawn for each}
  \seqnote{8.8}{-4.25}{8.6}{each node's 60 ms window closes; all three saw all 3 bids, so all three\\ work out the same winner by themselves: node 3, with the lowest expected wait}
  \seqself{13.2}{-5.2}{73 ms: wins,\\ plans on the\\ worker pool}
  \seqnote{8.8}{-6.4}{8.6}{73 to 162 ms: the result goes through the shared log (Section~\ref{sec:sharedlog}):\\ node 3 saves its batch, sends it, and applies it once all members passed that round}
  \seqself{13.2}{-7.3}{162 ms: applied\\ at index 1,246}
  \seqrep{13.2}{0}{-8.3}{162.5 ms: Done(responder 34548, planned by node 3), one UDP datagram}
  \seqself{0}{-8.7}{187 ms: answered,\\ first attempt}
  \seqmsg{0}{4.4}{-9.6}{cancel}
  \seqmsg{0}{8.8}{-9.9}{}
  \node[mlabel] at (6.6,-9.9) {cancel};
  \seqnote{11.0}{-10.25}{4.0}{1,014 ms: node 3 proposes the release of the responder; applied at 1,060 ms}
\end{tikzpicture}
\caption{One report at 10,000 a second, from the logs (\commit{eba4d14}, \texttt{plan\_work} 0). Nodes 1 and 2 lost the bid and stood by; they say nothing to the client unless the winner goes quiet. The client gets one answer and cancels its calls to the other two nodes.}
\label{fig:seq-report}
\end{figure}

\subsection{Paths and messages}

\begin{table}[!htb]
\centering\small
\caption{The paths between programs.}
\label{tab:paths}
\begin{tabularx}{\linewidth}{@{}P{3.4cm}P{3.6cm}L@{}}
\toprule
Path & Transport & What it carries \\
\midrule
Client to all 3 nodes & tarpc over UDP, port 8001 & \code{submit(ClientRequest)}, three copies with one ID; a cancel once answered \\
Node to client & the same socket & one \code{ClientReply}; the first one wins \\
Node to node & tarpc over TCP, port 7001; the shared log has its own connection (\dlog{D251}) & log streams, bids, heartbeats, fault bids, catch-up \\
\texttt{ctl} and web page to a node & tarpc over TCP, port 8001 & status, silence, slowdown, records; passes the fault gate \\
\bottomrule
\end{tabularx}
\end{table}

\begin{table}[!htb]
\centering\small
\caption{The main messages and records (\texttt{src/proto.rs}). Every number is an integer.}
\label{tab:messages}
\begin{tabularx}{\linewidth}{@{}P{3.5cm}L@{}}
\toprule
Message & Fields and meaning \\
\midrule
\code{ClientRequest} & \code{id} (client, seq), \code{attempt}, \code{age\_ms}, \code{incident}. A retry keeps the ID and raises \code{attempt}. \code{age\_ms} is how old the report is on the client's own clock, so ``too old'' never compares two clocks (\dlog{D143}). \\
\code{Incident} & \code{owner}, \code{category}, \code{zone}, \code{location}, \code{severity} 1 to 5, \code{needs}, \code{reported\_at\_ms}, \code{ttl\_ms} (\dlog{D140}). \\
\code{ClientReply} & \code{id}, \code{from}, and a body: \code{Done(result)}, \code{Expired(result)} for a report too old to plan, or \code{Provisional(plan)} from a node alone (\dlog{D123}). \\
\code{Bid} & \code{request}, \code{attempt}, \code{score}. Sent in groups, one message per peer every 10 ms (\dlog{D247}, \dlog{D256}). \\
\code{FaultBid} & \code{round}, \code{ticket}, \code{ms}: one node's numbers in the fault election (\dlog{D158}). \\
\code{Heartbeat} & nothing but the sender and its view number; any answer means alive (\dlog{D79}). \\
\code{RoundsMsg::Stream} & \code{after}, \code{through}, \code{batches}: my batches after round \code{after}, and my stream has passed round \code{through}. Answered \code{Have\{have\}} after saving. \\
\code{RoundsMsg::Drop}, \code{DropCommit} & go on without a silent node, and where its stream ends (Section~\ref{sec:drop}). \\
\code{RoundsMsg::CatchUp}, \code{Join} & a returning node asks for the applied log, then asks to be added (Section~\ref{sec:rejoin}). \\
\code{RoundBatch} & \code{node}, \code{round}, \code{commands}, and an optional \code{Add} of a returning node. \\
\code{Command} & \code{RegisterResponders}, \code{CommitResult} or \code{ReleaseResponder}; each safe to apply twice. \\
\bottomrule
\end{tabularx}
\end{table}

\textbf{On the wire.} Between nodes, a TCP frame is a 4-byte length, one byte that says whether the body is compressed, and the body: tarpc's message in bincode. A client datagram holds one or more short messages, each a call ID, a deadline and the call, at most 1,400 bytes after compression. Messages to the same address wait up to 20 ms to share a datagram, and bids wait up to 10 ms, because on the lab's Wi-Fi the number of packets mattered more than their size (\dlog{D254}, \dlog{D255}, \dlog{D256}). Compression is LZ4's block format \cite{lz4}, written by us in about 100 lines (\texttt{node/squeeze.rs}).

\subsection{Timers}

\begin{table}[!htb]
\centering\small
\caption{Timers in the lab runs. Safety never depends on a timer; only how fast the system reacts does.}
\label{tab:timers}
\begin{tabularx}{\linewidth}{@{}P{4.6cm}rL@{}}
\toprule
Timer & Value & Meaning \\
\midrule
\texttt{bid\_window\_ms} & 60 ms & how long a node collects bids for a report, from when the report reached it \\
\texttt{round\_ms} & 5 ms & how often a node makes its next batch for the shared log \\
\texttt{peer\_heartbeat\_ms} & 50 ms & how often nodes check each other \\
\texttt{takeover\_missed\_heartbeats} & 20 & so 1 s of silence before a peer is suspected and dropped \\
takeover backstop & $\ge$ 1.5 s & a loser plans the report if the winner has said nothing (\dlog{D250}) \\
\texttt{client\_retry\_ms}, \texttt{\_max\_ms} & 2 s, 6 s & the client's waits: 2, 4, 6, 6, 6 and 6 s, each $\pm$20\% \\
\texttt{client\_max\_retries} & 5 & then the client returns an error \\
call deadline & 36 s & one per report: 1.2 times all the waits together (\dlog{D169}) \\
\texttt{propose\_timeout\_ms} & 3 s & how long a command may wait to be applied \\
\texttt{peer\_rpc\_timeout\_ms} & 500 ms & one small call between nodes \\
\texttt{degraded\_after\_ms} & 2 s & a node alone this long answers provisionally \\
\texttt{JOIN\_GRACE}, \texttt{BOOT\_GRACE} & 3 s, 3 s & time to send after being added; time to wait for others after a restart \\
\bottomrule
\end{tabularx}
\end{table}

The first choice for detection was 3 missed heartbeats, about 150 ms. On the lab's Wi-Fi this dropped healthy nodes under load, so we moved to 1 s (\dlog{D249}). A failure therefore costs at least 1 s, and we report it that way.

\subsection{Retries, duplicates and requests in flight}
\label{sec:retries}

\textbf{Retries.} The class notes say a client should try 5 times and then report an error. Our client sends a report, waits 2 s, and sends it again with the same ID, then waits 4 s, then 6 s, up to 5 retries, and then \code{submit} returns an error (\dlog{D52}). Each wait has $\pm$20\% random spread so that clients do not retry together. Test T5 checks this with all nodes off: the client gave up with \texttt{"attempts":6} (Section~\ref{sec:tests}).

During a long outage at high load, retries can flood the network that the returning node needs to catch up. We measured this: with a second node killed at 20,000 a second, clients sent about 30,000 retries a second and answers fell to zero (\dlog{D264}). So a client earns 0.2 of a retry for each first send, plus 200 a second, and may save up to 5,000. A short failure uses the savings and is not slowed; a long one cannot flood the air. A retry with nothing to spend is not sent, and its time still passes, so a report still gives up when it always did.

\textbf{Duplicates.} Every report has one ID, made of the client's number and a sequence number. The first result recorded for an ID wins on every node, and a later one gets the winner back (\dlog{D27}). So planning a report twice, by a takeover or by two nodes that both think they won, is wasted work but never a wrong answer. A node that already holds a result answers a retry from its own copy, and a command already decided never goes into the log again (\dlog{D37}, \dlog{D250}).

\textbf{Requests in flight at a failure.} A report in flight when a node fails is covered three ways. If the winner goes quiet, the next node in that report's ranking plans it once the winner is suspected, or at the backstop (Section~\ref{sec:takeover}). If the answer is lost, the client's retry reaches all nodes again, and any node that holds the result answers from its copy. If the winner's result was already in the log, any node answers with it. Section~\ref{sec:s3} measures how many reports took each path.

\textbf{The client's view.} Figure~\ref{fig:client-states} is the state machine of one report inside \code{MeshClient}.

\begin{figure}[!htb]
\centering
\begin{tikzpicture}[x=1cm,y=1cm]
  \node[sinit] (i) at (0,0) {};
  \node[state, fill=clientorange!10, text width=3.2cm] (sent) at (3.2,0) {\textbf{Sent}\\ one call per node, attempt $a$};
  \node[state, fill=clientorange!10, text width=3.0cm] (retry) at (3.2,-2.6) {\textbf{Retrying}\\ same ID, attempt $a+1$, if the budget allows};
  \node[state, fill=mwgreen!14, text width=3.6cm] (done) at (10.2,1.3) {\textbf{Done}: a plan on record};
  \node[state, fill=mwgreen!14, text width=3.6cm] (exp) at (10.2,0) {\textbf{Expired}: too old, nobody sent};
  \node[state, fill=notefill, text width=3.6cm] (prov) at (10.2,-1.3) {\textbf{Provisional}: a lone node's plan};
  \node[state, fill=badred!10, text width=3.6cm] (gave) at (10.2,-2.6) {\textbf{Gave up}: \code{submit} returns an error};
  \draw[flow] (i) -- node[elabel, above]{new ID} (sent);
  \draw[flow] (sent.east) -- node[elabel, above, sloped, pos=0.55]{first \code{Done}} (done.west);
  \draw[flow] (sent.east) -- node[elabel, above, pos=0.55]{first \code{Expired}} (exp.west);
  \draw[flow] (sent.east) -- node[elabel, below, sloped, pos=0.55]{first \code{Provisional}} (prov.west);
  \draw[flow] ([xshift=-6pt]sent.south) -- node[elabel, left]{no answer in the wait} ([xshift=-6pt]retry.north);
  \draw[flow] ([xshift=6pt]retry.north) -- node[elabel, right]{resend} ([xshift=6pt]sent.south);
  \draw[flow] (retry) -- node[elabel, above]{5 retries used} (gave);
\end{tikzpicture}
\caption{One report as the client middleware sees it. On any answer the client cancels its calls still open on the other nodes. An answer to an earlier attempt still counts.}
\label{fig:client-states}
\end{figure}

\subsection{Choice of RPC}
\label{sec:rpc}

The class notes of Oct 4 ask us not to go down to raw UDP but to use a ready-made RPC or RMI system and say why, and the notes of Oct 28 point at the table in Chapter 2 of Java RMI, CORBA and Sun RPC. We chose RPC, not RMI, and the library \texttt{tarpc} \cite{tarpc} as the system (\dlog{D119}, \dlog{D176}). Our service is procedures over replicated state: nothing in it has an object identity that a client would name or hold, which is what RMI adds. Rust rules out using Java RMI or CORBA as they are, so we compared tarpc with them on the same rows (Table~\ref{tab:rpc}).

\begin{table}[!htb]
\centering\small
\caption{tarpc against the systems in the slides (\dlog{D176}).}
\label{tab:rpc}
\begin{tabularx}{\linewidth}{@{}P{2.6cm}LLLL@{}}
\toprule
 & Java RMI & CORBA & Sun RPC & tarpc (ours) \\
\midrule
Style & RMI & RMI & RPC & RPC \\
Languages & Java only & many & many & Rust only \\
Interface & remote interface & CORBA IDL & Sun XDR & a Rust trait with \texttt{\#[tarpc::service]} \\
Marshalling & Java serialization & CDR & XDR & serde with bincode \\
Semantics alone & at-most-once & at-most-once & at-least-once & maybe: it never resends \\
\bottomrule
\end{tabularx}
\end{table}

tarpc on its own gives ``maybe'' semantics, so at-most-once is our work (\dlog{D177}). From client to cloud, the effect of a report is recorded at most once across all replicas: the client resends with the same ID, the results table filters duplicates by ID, and any node re-sends the answer from its copy. Peer calls are ``maybe'': a lost bid or heartbeat is simply covered by the next one. The one labelled exception is a provisional answer from a node alone, which can differ from what is recorded later (Section~\ref{sec:degraded}).

\textbf{Why UDP for reports.} The note ``Request reply protocol + UDP'' can be read as ``use UDP'' or as ``UDP is the level not to go down to''. We did both: reports travel over UDP, but our UDP adapter holds no RPC logic. Call IDs, deadlines and matching replies to calls stay tarpc's, and retries stay in the middleware (\dlog{D178}). UDP saves the connection set-up per client and node, the pings and reconnects, and acknowledgements that a reply already gives. It costs a full wait when an answer is lost, a larger kernel buffer and a firewall rule.
